\documentclass{article} %210 mm × 297 mm
\usepackage[utf8]{inputenc}
\usepackage[a4paper,left=1.5cm, right=1.5cm, top=2cm, bottom=2cm]{geometry}
\usepackage{graphicx}
%\usepackage{blindtext}
\usepackage{multirow}
\usepackage{mhchem}
\usepackage{url}
\usepackage{subfigure}
\usepackage{amsfonts,latexsym} % para tener disponibilidad de diversos simbolos
\usepackage{enumerate}
%\usepackage{booktabs}
\usepackage{float}
%\usepackage{threeparttable}
\usepackage{array,colortbl}
%\usepackage{ifpdf}
%\usepackage{rotating}
\usepackage{cite}
\usepackage{stfloats}
\usepackage{url}
\usepackage{listings}
\usepackage{fancyhdr}

\usepackage{color}
%\usepackage{hyperref}
%\usepackage{wrapfig}
\usepackage{array}
%\usepackage{multirow}
%\usepackage{adjustbox}
%\usepackage{nccmath}
\usepackage[dvipsnames]{xcolor}

 
\newcommand{\titulo}{Präsenzübungsblatt 8; Statistik}
\newcommand{\nombre}{Lil'Jana}
\newcommand{\FCE}{\textbf{SoSe 2025}}
\newcommand{\dat}{\textbf{Tut 03.06.2025}}
\newcommand{\gre}[1]{\textcolor{green}{#1}}
\newcommand{\purp}[1]{\textcolor{purple}{#1}}
\newcommand{\cyan}[1]{\textcolor{cyan}{#1}}
\newcommand{\yel}[1]{\textcolor{yellow}{#1}}
\newcommand{\blu}[1]{\textcolor{blue}{#1}}
\newcommand{\red}[1]{\textcolor{red}{#1}}
\newcommand{\p}[0]{\mathbb{P}}
\newcommand{\E}[1]{\mathbb{E}[#1]}

\usepackage{fancyhdr}
\pagestyle{fancy}
\fancyhead{}
\fancyfoot{}
\fancyhead[RO]{\small\dat}
\fancyhead[LO]{\small\FCE}
\fancyfoot[RO,LE] {\thepage}

\setcounter{page}{1} %Número de la página, la editorial lo cambiará

\newcommand{\kl}[1]{\{ #1 \}}

\usepackage[onehalfspacing]{setspace}
\begin{document}


\begin{tabular}{p{12cm}}
{{\Large\bf\titulo}}\\
\vspace{0.2cm}\\
 {\large\nombre}\\
\end{tabular}
\vspace{0.2cm}\\

\section*{9.1}
\subsection*{(a)}
\subsubsection*{Aufgabe bedeutet:}
\[
X \sim Poi(\lambda x) \lambda x > 0 \quad x \in \mathbb{N}
\]
\[
\p(X=n) = \frac{\lambda x^n}{n!} e^{-\lambda x}
\]
\[
Z \sim Geo(p) \quad p \in (0,1), \quad Z \in \mathbb{N}
\]
\[
\p(Z=n) = (1-0)^{n-1} p \quad (n \geq 1)
\]
\subsubsection*{In den Folien steht:}
Wenn $X$ und $Y$ unabh.: 
\[ 
\E{X \cdot Y} = \E{X} \cdot \E{Y} 
\]
und für $z \neq 0$ gilt: 
\[
\p (XY = z) \sum_{y \neq 0} \p(XY = z \mid Y = y) \cdot \p (Y = y) \overset{*}{=}  \sum_{y \neq 0} \p (X = \frac{z}{y})  \cdot \p (Y = y)
\]
* nur weil $X$ und $Y$ unabhängig sind! \\
Außerdem: 
\[
Var(X) = \E{(X - \E{X^2})} = \sum_{x \in X(\Omega)} (x - \E{X})^2 \cdot \p (X=x)  
\]
Und: Poisson-verteilte Zufallsvariable $X$ mit Parameter $\lambda > 0$:
\[
\E{X} = \lambda = \Var(X)
\]
Und: geometrisch verteilte Zufallsvariable $X$ mit Erfolgswahrscheinlichkeit $p$: 
\[
\E{X} = \frac{1}{p} \quad und \quad Var(X) = \frac{1-p}{p^2}
\]
\subsubsection*{An die Tafel hat sie noch geschrieben:}
\[
Var(X+Y) = Var(X) + Var(Y) \underbrace{ + Cov(X,Y)}_{=0\ wenn\ unabh.}
\]
\[
Var(aX + b) = a^2 Var(X)
\]
\subsubsection*{Also rechnet man (meine Rechnung)}
 \[
 \E{X\cdot Z} = \E{X} \cdot \E{Z} = \lambda_X \cdot \frac{1}{p} \quad \textcolor{orange}{\text{korrekt}}
 \]
 und 
 \[
 Var(Y + Z + 5) = Var(Y) + Var(Z) + Var(5)  = \textcolor{red}{\E{(Y - \E{Y^2})}} + \frac{1-p}{p^2} + \textcolor{red}{Var(5)} \quad \textcolor{orange}{\text{nicht ganz}}
 \]
 \textcolor{orange}{
 \[
 Var(Y) + Var(Z) + 0 = \lambda_Y + \frac{1-p}{p^2} 
 \]
 Und die letzte: 
 \[
 Var(-2\cdot Z + X) = Var(-2\cdot Z) + Var(X) = \underbrace{4}_{(-2)^2} \cdot Var(Z) + Var(X) = 4 \cdot \frac{1-p}{p^2}  + \lambda_X
 \]
 }
 \subsection*{b}
 \subsubsection*{1.}
 Von der Tutorin:  \\
 z.z. \(X_1^2\)und $X_2^2$ unabh. $\rightarrow$ z.z. \(\p(X_1^2 \leq S, X_2^2 \leq t) = \p(X_1^2 \leq S) \cdot \p(x_2^2 \leq t)\)
 \[
\underbrace{\p(X_1^2 \leq s, X_2^2 \leq t)}_* = \p(X_1 \leq \sqrt{s}, X_2 \leq \sqrt{t}) \overset{**}{=} \p(X_1 \leq \sqrt{s}) \cdot \p(X_2 \leq \sqrt{t}) = \p(X_1^2 \leq s) \cdot \p(x_2^2 \leq t)
 \]
 * $X_1^2 \in [0,s]$ (da $X_1^2 \geq 0$) $\Rightarrow x_1 \in [-\sqrt{s}, \sqrt{s}]$   \\
 ** weil unabh. 
 
 
 \subsubsection*{2.}
 \[
 Var(X \cdot Y ) = \E{(XY)^2} - (\E{XY}
 \]
 

\end{document}